\begin{table}[htbp]
\centering
\renewcommand{\arraystretch}{1.2}
\begin{tabular}{|p{2.2cm}|p{4.2cm}|p{4.2cm}|p{3.2cm}|}
\hline
\textbf{Frequency} & \textbf{Ailment / Outcome} & \textbf{Target Metabolic Cycle} & \textbf{Study} \\
\hline
40\,Hz & Tau pathology $\downarrow$, motor function $\uparrow$ & Neural gamma entrainment, proteostasis & Iaccarino et al., 2016; Martorell et al., 2019 \\
\hline
0.8\,Hz & Sleep memory $\uparrow$ & Cortical slow oscillation & Ngo et al., 2013 \\
\hline
30--50\,Hz & Vagus stimulation, inflammation $\downarrow$ & Afferent vagal tone modulation & Bonaz et al., 2016 \\
\hline
62--125\,Hz & Mitochondrial ROS / number modulation & Cellular EM field resonance & Krajnak et al., 2020 \\
\hline
100\,Hz & Neurite growth $\uparrow$ & Cytoskeletal reorganization & Mo et al., 2022 \\
\hline
\end{tabular}
\end{table}

This survey highlights a recurring pattern in the published literature: discrete frequency bands are associated with measurable physiological effects that map, at least provisionally, onto specific oscillatory or metabolic processes. The examples above are not presented as a unified mechanism, but as a compact reference set for the broader RBM hypothesis that biological systems may be selectively responsive to frequency-matched stimulation.
